% Khung bài tập tín hiệu 4 mức, vẽ lại từ hình mẫu.
% Dùng trong tài liệu đã nạp gói tikz bằng lệnh \input{...}.
\begin{tikzpicture}[
    x=1.5cm,
    y=0.6cm,
    font=\small\bfseries\boldmath,
    line width=1.1pt,
    line cap=round
]
    \begin{scope}[draw=black, text=black]
        % Sáu vạch chia tạo thành năm khoảng ký hiệu, mỗi khoảng hai bit.
        \foreach \boundary in {0,...,5} {
            \draw (\boundary, -3.6) -- (\boundary, 3.6);
        }

        % Đường ngang ở giữa khung.
        \draw (0, 0) -- (5.85, 0);

        \node[anchor=east] at (-0.45, 3.6) {$S_t=$};
        \node[anchor=east, align=left] at (-0.45, 0)
            {Tín hiệu 4\\mức};

        % Giữ nguyên thứ tự các cặp bit trong ảnh.
        \foreach \position/\bits in {
            0.5/{1\,0},
            1.5/{0\,1},
            2.5/{1\,1},
            3.5/{0\,0},
            4.5/{0\,1}
        } {
            \node at (\position, 3.05) {$\bits$};
        }

        % Bốn mức biên độ; hoàn chỉnh các nhãn bị cắt ở mép ảnh.
        \foreach \level in {3, 1, -1, -3} {
            \node[anchor=west] at (5.65, \level) {$\level$};
        }
    \end{scope}
\end{tikzpicture}
